\section{What remains open, and what was not checked}\label{sec:open}

\subsection{Open questions}

\begin{enumerate}[label=\arabic*.]
\item \emph{The construction and the counterexample.} Whether the constructed threefold exists with the stated properties, and so contradicts Theorem~2.2 of \cite{CDP20}, needs specialist verification. The record's two \textsf{OPEN} rows are exactly this: independent replication of the cyclic invariant lattices and of the cotangent calculation. Engel's exposition \cite{EngelS6} is a further independent source, and its comparison with the record has not been made.
\item \emph{The published step.} Whether the proof in \cite{CDP20} uses the implication of Section~\ref{sec:correction} in the unrestricted form that the test family refutes needs a reading of the paper, which was not made here. By Corollary~\ref{cor:normal}, the implication can fail only on fibres that are singular along a curve.
\item \emph{The global question left by the note.} Whether the global hypotheses of \cite{CDP20} force $\delta_A$ to be injective on such fibres is open. The test family shows that no local argument can give this: there the target of $\delta_A$ vanishes (Proposition~\ref{prop:sharp}).
\item \emph{Relative two-forms at the cusp.} The torsion of $\Omega^2_{X/B}$ on the record's charts is computed in Proposition~\ref{prop:torsion}(d): along the double curves it is the structure sheaf of the curve, and at a triple point the canonical module of the three branches. Its global form along the cusp fibre, and its effect on the direct images $R^qf_*\Omega^p_X$ through which the manuscript's Remark~1.3 would compute the Hodge numbers, were not examined. The record computes $h^{1,2}=1$ by another route (its ledger row~17).
\item \emph{The deformation of S6-2.} The record does not claim that different values of the parameter give non-isomorphic complex manifolds, or that the family exhausts the deformations (key-advances reader, item~2).
\item \emph{The two glue codes.} The Erdős--Straus workbench's order-$72$ glue code for $A_5^4D_4$ and the one used in S6-4 have not been compared (Section~\ref{sec:overlaps}).
\end{enumerate}

\subsection{What this reader did not check}

\begin{itemize}
\item \emph{The construction.} Sections~6--7 of the monograph (the period-family descent, the fillings, the gluing, the nearby-cycle and Leray comparisons) were not read. The topology and sphere recognition, the algebraic dimension, and the Hodge, Bott--Chern and Aeppli computations were not read either. Every statement about them in this reader is the record's.
\item \emph{The monograph beyond the parts named.} Of its 154 pages, Part~I, the ledger, Theorem~13.6 and §§13.11 and~14 were read. The rest was mapped from its table of contents.
\item \emph{The later papers and the archive.} The two papers of the topology supplement (\fn{15\_}, \fn{16\_}) were read only through their guide and their front matter. The proofs of S6-1 to S6-4 are in the archive \fn{28\_}, which was not downloaded.
\item \emph{The conversation and workflow files} \fn{07\_}--\fn{13\_} were not used.
\item \emph{The literature.} The paper \cite{CDP20} was not re-read; its statements are cited from the record and the note. No novelty search was made for the additions of Section~\ref{sec:correction}. The normal case of Corollary~\ref{cor:normal} is the source manuscript's Remark~10.4, and the general facts used are cited where they are used.
\end{itemize}
